% New v7 material; does not modify the v6 bundle.
% Existing bibliography keys used: EM, Evans, GKM, Hattori, PapadimaSuciu.
% Existing labels used: snc:main, trop:petersen-unimodular,
% trop:farey-condition, trop:branch-denominator.

\section{Petersen core の作用座標を実現する面積の選択}
\label{per:section}

Petersen graph と標準 fan の同型だけでは，実際の作用周期の整合性は決まらない．
一方，foliation の相対 isotopy class を保ちながら代表を選べる場合には，
五直線補集合の一つの core で，この整合性を具体的に実現できる．
必要な操作は，各 puncture に接する critical lobe の面積を，その puncture の residue に合わせることである．
これにより十五の全 face が unimodular になり，三価 branch での格子も飽和する．

contact link の作用積分と symplectization による cone construction は
Evans--Mauri\cite[Remarks 2.11--2.12, Lemma 2.16]{EM}に基づく．
標準 Petersen fan の格子は，命題\ref{trop:petersen-unimodular}および
Gathmann--Kerber--Markwig\cite[Construction 3.6, Theorem 3.7]{GKM}の unit-edge lattice を用いる．
以下の構成で追加するのは，実際の作用積分とこの格子を合わせる面積条件と明示的な face maps である．

\subsection{適合性の意味と面積指定の補題}

正則 torus の整数一次ホモロジーを $L_e=H_1(T_e;\ZZ)$ とすると，
作用一形式 $\dd I_\gamma$ により $L_e$ は底の整数余接線格子となる．
従って整数接線格子は $L_e^*=\operatorname{Hom}(L_e,\ZZ)$ である．
標準 fan の ambient tangent lattice $N$ と $L_e$ 自体を同一視しない．

\begin{definition}[作用座標に適合する cone chart]
\label{per:chart-definition}
Liouville cone 型 fibration の底の core 頂点を $o$ とする．
原点を持つ二次元 integral fan $\mathcal F\subset N_\RR$ への
\emph{作用座標に適合する cone chart} とは，$o$ の近傍から $0$ の近傍への
stratum を保つ同相写像 $F$ であって，次を満たすものである．
各正則 face $e$ 上で，$F$ は作用座標の整数線形変換であり，
その微分は $L_e^*$ を $N\cap\operatorname{span}_\RR F(e)$ へ同型に写す．
さらに，branch ray の接線格子は，各 incident face の接線格子の飽和部分格子である．
\end{definition}

primitive を固定すると作用座標は core に向かって零に収束するので，
$F(o)=0$ という正規化の下で translation は零となる．
この定義は全空間の任意の homology class に ambient linear function を割り当てることを要求しない．
その追加条件との違いは後で明示する．

\begin{lemma}[四つの critical lobes の面積指定]
\label{per:lobe-areas}
面積形式 $\omega_C$ を持つ球面 $C$ に四つの異なる marked points を指定し，
それらの $2+2$ 分割を固定する．
この分割を実現する，二つの simple saddles を持つ marked Morse foliation の
相対 isotopy class を一つ固定する．
四つの正数 $s_1,s_2,s_3,s_4$ が
\[
 s_1+s_2+s_3+s_4<\int_C\omega_C
\]
を満たすとする．十分小さい marked disks を固定したまま，この相対 isotopy class の中で，
四つの external Reeb edges に対応する critical basins の面積をそれぞれ $s_j$ とする代表を選べる．
ここで critical basin とは，puncture から external edge に沿って拡大する regular disks の増大極限である．
その境界が critical figure-eight の単純な一方の loop であることは要求しない．
二つの critical levels の間の regular annulus の面積は
$\int_C\omega_C-\sum_js_j$ となる．
\end{lemma}

\begin{proof}
初めに，指定した class の代表を一つ取る．
二つの figure-eight critical levels は，球面を四つの open critical basins と一つの open annulus に分ける．
critical levels 自体の面積は零である．
各 marked disk は対応する lobe の内部に十分小さく取り，その面積を $s_j$ より小さくする．

元の面積形式に正の smooth density を掛け，marked disks の近傍では density を一に保つ．
それ以外ではまず density を十分小さくし，各 open region 内の非負 bump functions を加えることで，
四つの critical basins と annulus の積分をそれぞれ指定した五つの正数にできる．
こうして得る面積形式を $\widehat\omega_C$ とする．
$\int_C\widehat\omega_C=\int_C\omega_C$ であり，両形式は marked disks の近傍で一致する．

差は exact である．その primitive は，各 disk 上では closed なので，
局所 primitive の cutoff を引くことにより marked disks の近傍で零にできる．
$\omega_C$ と $\widehat\omega_C$ の凸結合に Moser の議論を適用すると，
disks を点ごとに固定する isotopy を得る．
この isotopy によって元の foliation を運べば，$\omega_C$ に関して指定した面積を持つ代表となる．
Morse 型および相対 isotopy class は保存される．
\end{proof}

この補題の可動対象は foliation の代表である．
既に固定された embedded critical leaves を点ごとに保ったまま，その囲む面積を変更する主張ではない．

\subsection{五直線に対する residue 正規化}

$D=L_0\cup\cdots\cup L_4\subset\PP^2$ を一般位置の五直線とし，
$w_a>0$ および
\begin{equation}\label{per:positive-weights}
 [\omega]=2\pi\sum_{a=0}^4w_a[L_a]
\end{equation}
を固定する．各 $L_a$ は面積 $2\pi\sum_bw_b$ を持つ．
crossing charts と normal form の準備後，その時点の divisor 上の面積形式を用いて，
以下の marked foliations の代表を選ぶ．この順序により，後の normal coupling は divisor 上の形式を変えない．

添字は $\ZZ/5\ZZ$ とする．$L_a$ の二つの saddle vertices を $A_a,B_a$ とし，
$A_a$ 側の punctures を $a-1,a+1$，$B_a$ 側を $a-2,a+2$ とする．
link の辺は
\[
 A_aB_a,\qquad A_aA_{a+1},\qquad B_aB_{a+2}\quad(a\in\ZZ/5\ZZ)
\]
であり，これは Petersen graph である．
同じ組合せ型の任意の marked relative isotopy classes を出発点にしてよい．

補題\ref{per:lobe-areas}を各 $L_a$ に適用し，puncture $L_a\cap L_b$ に接する
critical basin $D_{a;b}$ の面積を
\begin{equation}\label{per:normalization}
 \operatorname{area}(D_{a;b})=2\pi w_b\quad(b\ne a)
\end{equation}
とする．その四つの和は全体より $2\pi w_a$ だけ小さい．従って，内部 annulus $E_a$ の面積は
\begin{equation}\label{per:annulus-one}
 \operatorname{area}(E_a)=2\pi w_a.
\end{equation}
十個の crossing disks は先に十分小さく選び，面積を各 target より小さくしておく．
各 component の相対 isotopy はその disks を固定するので，相互の crossing data に干渉しない．

\begin{lemma}[critical crossing periods の消去]
\label{per:zero-periods}
比例曲率 collar で $\lambda=(m_a-w_a)\alpha_a$ とする．
$L_a\cap L_b$ の product chart における正の meridian を $\mu_b$ と書き，
その cycle を対応する regular sheet に沿って critical lobe まで運ぶ．
正規化 \eqref{per:normalization} の下で，critical ray 上のその作用積分は零である．
normal meridian $\mu_a$ の作用積分は $m_a-w_a<0$ であり，
ここで $m_a$ は normal radial moment を表す．
\end{lemma}

\begin{proof}
cycle と radial moment の記号を区別して計算する．
product chart の puncture-compatible framing に沿って degree-one lift を取る．
critical lobe $D_{a;b}$ に対する connection holonomy の実数 lift は，Stokes の定理から
\begin{equation}\label{per:holonomy-lobe}
 h_{a;b}=
 \frac{w_b-\operatorname{area}(D_{a;b})/(2\pi)}{w_a}
\end{equation}
となる．実際，$L_a$ 内の小さい circle の disk area を $2\pi I$ とすれば，
normal collar に入った所での holonomy は $(w_b-I)/w_a$ であり，
残る領域で $\dd\alpha_a=-\omega_{L_a}/w_a$ を積分すればよい．
\eqref{per:normalization} により $h_{a;b}=0$ となり，
$I_{\mu_b}=(m_a-w_a)h_{a;b}=0$ を得る．
normal cycle は $I_{\mu_a}=m_a-w_a$ である．
これらの恒等式は Liouville flow の下で同じ倍率で拡大・縮小されるので，critical ray 全体で成り立つ．
相対構成の graph-period correction は torus tangent を消す basic forms によるため，
これらの fiberwise action periods を変えない．
\end{proof}

\subsection{十五の face maps と core chart}

$\mu_a$ は以後 meridian cycle だけを表す．
五直線補集合 $W_5$ の整数一次ホモロジーは
\begin{equation}\label{per:global-meridians}
 H:=H_1(W_5;\ZZ)
  =\left(\bigoplus_{a=0}^4\ZZ\mu_a\right)
     /\ZZ(\mu_0+\cdots+\mu_4)
\end{equation}
である\cite{Hattori,PapadimaSuciu}．
crossing torus 上の二つの meridians は整数基底をなす．
内部 annulus の torus 上では，
\[
 \alpha_a=\mu_a,\qquad
 \beta_a=\mu_{a-1}+\mu_{a+1}
\]
が整数基底となる．$A_a$ 側の pants block で後者は二つの lobe cycles の和である．
$B_a$ 側では，同じ annular cycle は
\begin{equation}\label{per:normal-euler-framing}
 \beta_a=-\mu_a-(\mu_{a-2}+\mu_{a+2})
\end{equation}
となる．これは normal bundle の Euler 数 $L_a^2=1$ に対応する framing relation である．

ここで \eqref{per:global-meridians} を作用積分に直接代入してはいない．
$\lambda$ は closed ではないので，全空間で homologous な曲線の積分が一致するとは限らない．
まず実際の annular torus の二つの端から運んだ cycles を比較する．
この torus の $H_1$ から $H$ への写像は，$\mu_a$ と $\mu_{a-1}+\mu_{a+1}$ の線形独立性から injective である．
従って \eqref{per:global-meridians} は，この torus 自体の中で
\eqref{per:normal-euler-framing} を与える．その後に，torus 上で closed な $\lambda$ の積分を取る．

標準 Petersen fan の rays を $r_{A_a},r_{B_a}$ とし，前節のラベル付けを使う．
例えば $A_a\leftrightarrow\{2a,2a+1\}$，$B_a\leftrightarrow\{2a+2,2a+4\}$ を
$\ZZ/5\ZZ$ 上で取り，最後にラベルを $1,\ldots,5$ に読み替えればよい．

\begin{theorem}[一つの Petersen core における作用座標の実現]
\label{per:local-realization}
\eqref{per:positive-weights} を満たす任意の正の weights に対し，
Petersen 型の指定された marked relative isotopy classes の中から代表を選び，
相対 SNC 構成による $W_5$ の core neighborhood に
定義\ref{per:chart-definition}の作用座標に適合する標準 Petersen cone chart を与えられる．
全ての branch denominator は $q=1$ であり，十五の face maps は unimodular である．
この主張は一つの core の近傍についての存在を述べる．
\end{theorem}

\begin{proof}
面積条件 \eqref{per:normalization} を実現し，補題\ref{per:zero-periods}を用いる．
crossing face $e=vw$ が $L_a,L_b$ を結ぶとき，その face 上で
\begin{equation}\label{per:crossing-map}
 F=-I_{\mu_a}r_v-I_{\mu_b}r_w
\end{equation}
と定める．両作用が sector 全体で負であることを確認する．
component 側の external arm では，partial disk の面積が critical basin の面積より小さいため，
\eqref{per:holonomy-lobe}の holonomy は正である．normal action は負なので他方の meridian action も負となる．
transition cap は $I<w_b/2$ に取り，normal radius を一様に縮めれば
$I-w_b+m_a A_{ab}(I)<0$ である．product crossing でも $m_a<w_a$，$m_b<w_b$ とすれば両係数は負となる．
$v$ 側の ray では $I_{\mu_b}=0$，
$w$ 側では $I_{\mu_a}=0$ となる．
$\mu_a,\mu_b$ は torus homology の基底なので，係数への変換行列は $-I_2$ である．

内部 matching face $A_aB_a$ では
\begin{equation}\label{per:internal-map}
 F=(-I_{\mu_a}-I_{\beta_a})r_{A_a}+I_{\beta_a}r_{B_a}
\end{equation}
と定める．$A_a$ 側では $I_{\beta_a}=0$ である．
$B_a$ 側では \eqref{per:normal-euler-framing} と二つの lobe periods の消去から
$I_{\beta_a}=-I_{\mu_a}$ となる．
内部 annulus では Stokes の定理により holonomy が $0$ から $-1$ へ単調に変わるので，
二つの係数はいずれも正である．
$(I_{\mu_a},I_{\beta_a})$ から係数への変換行列は
\[
 \begin{pmatrix}-1&-1\\0&1\end{pmatrix}
\]
であり，determinant は $-1$ である．

三つの incident faces から $v\in\{A_a,B_a\}$ へ近づくと，いずれも
\begin{equation}\label{per:common-ray}
 F=-I_{\mu_a}r_v
\end{equation}
に一致する．従って各 map は branch 上で貼り合う．
$\mu_a$ は primitive で，作用座標における ray は各 face で primitive integer vector を持つため，
その branch lattice は飽和する．
具体的に crossing face では $(-1,0)$，内部 face では $(-1,0)$ または $(-1,1)$ となる．

各 sector が所期の fan sector の近傍へ一対一に写ることも確認する．
contact link の regular edge を $t$ で parametrise し，作用 vector を $a(t)$ とすると，
symplectization の作用 map は $(r,t)\mapsto r a(t)$ である．
contact form の作用 vector は零にならず\cite[Remark 2.11]{EM}，
正則 torus fibration の作用座標は局所座標である．従って
$\det(a(t),a'(t))\ne0$ である．
上で示した正の sector 内では，projective action ratio は単調であり，両境界 rays の間を一度だけ走る．
従って map は sector の近傍の同相写像となる．
Liouville flow は作用積分を同じ倍率で縮めるため，compact な link 全体で core への収束は一様である．
有限個の sector maps を貼り，必要なら近傍を縮めれば，fan の原点近傍への同相写像を得る．

最後に，標準 Petersen fan の全 cones はそれぞれの rank-two span で unimodular である．
\eqref{per:crossing-map} と \eqref{per:internal-map} の整数行列も unimodular なので，
各 face の action cotangent lattice と fan tangent lattice の dual が正しく一致する．
\end{proof}

\begin{lemma}[branch star における二つの整数延長]
\label{per:branch-duality}
前定理の一つの branch $v$ で，crossing neighbors を $x,y$，matching neighbor を $z$ とする．
rank-three star lattice の基底を $(r_v,r_x,r_y)$ と取れば
$r_z=r_v-r_x-r_y$ である．
$H_1(Q_v)=\ZZ\alpha\oplus\ZZ b_1\oplus\ZZ b_2$ とし，二つの crossing faces の cycle lattices を
$\langle\alpha,b_1\rangle$，$\langle\alpha,b_2\rangle$，matching face のものを
$\langle\alpha,b_1+b_2\rangle$ とする．
各 face の作用積分を保つ，$H_1(Q_v)$ から star cotangent lattice への同型は，
この基底を固定すればちょうど二つあり，その一つは
\[
 \alpha\longmapsto(-1,0,0),\qquad
 b_1\longmapsto(0,-1,0),\qquad
 b_2\longmapsto(0,1,-1)
\]
で与えられる．
\end{lemma}

\begin{proof}
この三本の rays は saturated rank-three lattice を生成する．
例えば $v=15$ の場合は $r_v,r_{23},r_{24}$ の三次 minor が $\pm1$ であり，
他の場合は五つの labels の置換で従う．
face restrictions により，延長は
\[
 \alpha=(-1,0,0),\qquad b_1=(0,-1,t),\qquad b_2=(0,s,-1)
\]
という形に限る．matching ray 上で $b_1+b_2$ の値は一なので $s+t=1$ である．
この三行の determinant の絶対値は
$1-st=s^2-s+1$ となる．これが一であるのは $s=0,1$ の二つだけである．
\end{proof}

\begin{remark}[この局所存在が与えないもの]
\label{per:scope}
定理\ref{per:local-realization}では，全空間の boundary contact model を標準 $M_{0,5}$ の filling に置き換えていない．
common-circle determinant と boundary homology の差はそのまま残る．
また，六個の core の間で共有される actual double curves 上の面積条件，
node patches，branch interiors および face transitions の同時整合は，この定理の範囲外である．
これらの追加条件は，節~\ref{sg8:section} の計算と定理~\ref{ga8:global} の相対接着で満たす．
\end{remark}

\subsection{固定された面積に対する奇数逆数条件}

\begin{proposition}[Farey 条件と branch saturation の結合]
\label{per:odd-reciprocal}
同一成分内の regular annulus $E$ の両端で branch tangent lattice が飽和し，
$E$ に対応する face が一つの unimodular cone となるとする．
その normal residue を $w>0$ とすると
\begin{equation}\label{per:odd-area-spectrum}
 \frac{\operatorname{area}(E)}{2\pi w}=\frac1m,
 \qquad m\text{ は正の奇整数}
\end{equation}
である．両端の reduced denominators $q,q'$ は互いに素な奇数で，$m=qq'$ を満たす．

逆に，任意の正の奇整数 $m$ に対して，branch numerator の primitivity と
一つの annulus の signed Farey equation からなる算術条件には解がある．
この逆向きの主張は算術方程式についてだけであり，指定された全ての geometric periods の実現を主張しない．
\end{proposition}

\begin{proof}
命題\ref{trop:branch-denominator}により $q,q'$ は奇数である．
命題\ref{trop:farey-condition}の determinant が $\pm1$ なので $\gcd(q,q')=1$ であり，
同じ命題の面積公式から \eqref{per:odd-area-spectrum} を得る．

逆向きには $m=qq'$ を互いに素な奇数の積に分ける．
内部 annulus の framing を \eqref{per:normal-euler-framing} と同じ convention に取り，
両端の internal numerators を $p,p'$ とする．正の面積 $1/(qq')$ に対応する式は
\[
 qq'+q'p+qp'=1.
\]
B\'ezout により整数解があり，式を $q$ および $q'$ で見ると
$\gcd(p,q)=\gcd(p',q')=1$ が従う．
奇数 $q$ では任意の unit $p$ を二つの units の和にできる．
例えば $a\equiv2^{-1}p\pmod q$，$b=p-a$ とすればよい．
$q'$ 側も同様である．$q=1$ の場合も $\gcd(0,1)=1$ の convention で成立する．
従って二つの branch の三 numerator に必要な primitivity を同時に満たせる．
\end{proof}

\begin{lemma}[branch numerator の分解数]
\label{per:unit-split-count}
$q\ge1$ と $\gcd(p,q)=1$ に対し，
$\gcd(a,q)=\gcd(p-a,q)=1$ を満たす residue class $a\pmod q$ の個数は
\[
 \prod_{\ell^e\parallel q}\ell^{e-1}(\ell-2)
\]
である．空積は一とする．従って解が存在するための必要十分条件は $q$ が奇数であることである．
また $m$ の互いに素な正整数による ordered factorization $m=qq'$ の個数は
$2^{\omega(m)}$ である．ここで $\omega(m)$ は異なる素因数の個数である．
\end{lemma}

\begin{proof}
素数 $\ell$ に対して除外されるのは $a\equiv0,p\pmod\ell$ の異なる二つの類である．
従って $\ell^e$ 上では $\ell^{e-1}(\ell-2)$ 個であり，Chinese remainder theorem で積を取ればよい．
互いに素な factorization では各 prime power 全体を二つの因子のいずれか一方に割り当てる．
\end{proof}

\eqref{per:odd-area-spectrum} は，固定された weights と critical leaves に対する離散的な必要条件である．
定理\ref{per:local-realization}は，そのうち $m=1$ を正の lobe areas によって幾何的に実現している．

\subsection{四直線の \texorpdfstring{$K_4$}{K4} core と有向四角形}
\label{per:k4-section}

三 punctures の場合にも，critical basin の面積を選ぶことで作用座標を整合させられる．
ここでは標準 tropical plane の格子
\[
 N_4=\ZZ^4/\ZZ(1,1,1,1),\qquad r_i=[e_i]
\]
を用いる．$\sum_i r_i=0$ であり，六つの二次元 cones は
$\RR_{\geq0}r_i+\RR_{\geq0}r_j$ $(i<j)$ である．
これは $\mathcal B(U_{3,4})$ の標準 fan で，その link は $K_4$ である．

\begin{proposition}[四直線 core の作用座標]
\label{per:k4-realization}
$W_4=\PP^2\setminus\bigcup_{i=1}^4L_i$ を一般位置の四直線の補集合とし，
$\omega$ を $[\omega]=2\pi\sum_iw_i[L_i]$，$w_i>0$ を満たす Kähler form とする．
各 $L_i$ の三 punctures 上で指定した marked Morse foliation の相対 isotopy class の中で
代表を選び，定理~\ref{snc:main} の相対構成による core に，
定義~\ref{per:chart-definition} の作用座標に適合する $K_4$ cone chart を構成できる．
branch denominator は全て $1$ である．

さらに，四頂点の有向 cycle $f$ を一つ選ぶと，この chart は整数格子同型
\begin{equation}\label{per:k4-lattice-map}
 \Phi:H_1(W_4;\ZZ)\longrightarrow N_4^*,\qquad
 \Phi(\mu_i)=e^*_{f^{-1}(i)}-e^*_i
\end{equation}
と整合するように取れる．ここで $\mu_i$ は正の meridian であり，
右辺は $\ZZ^4$ 上の係数和が零の linear form として $N_4$ に降りる．
整合性とは，各 regular face 上の torus cycle の作用積分が，
その image に対応する linear function の chart への pullback になることである．
\end{proposition}

\begin{proof}
crossing germs を準備した後の divisor 面積形式を用いる．
$L_i$ の一つの saddle の critical leaf は，球面を三つの puncture basins に分ける．
basin は external Reeb edge に沿う regular puncture disks の増大極限で定め，
その境界が一つの単純 loop であることを要求しない．
三つの面積の和は $2\pi W$，$W=\sum_jw_j$ である．
basin $D_{i;j}$ の面積を
\begin{equation}\label{per:k4-basin}
 \frac{\operatorname{area}(D_{i;j})}{2\pi}
 =w_j+w_i\,\mathbf1_{\{j=f(i)\}}\qquad(j\ne i)
\end{equation}
と指定する．これらは全て正で，和は $W$ に等しい．
補題~\ref{per:lobe-areas} の density と相対 Moser の証明を
三つの basins に適用すれば，十分小さい crossing disks を固定し，
指定された相対 isotopy class の中でこの面積を実現できる．
ここには内部 annulus がなく，三領域だけで全体の面積を尽くす．

puncture-compatible framing による Stokes の公式から，critical slope は
\[
 h_{ij}=\frac{w_j-\operatorname{area}(D_{i;j})/(2\pi)}{w_i}
       =-\mathbf1_{\{j=f(i)\}}.
\]
従って $i$ の ray 上では，$s=-I_{\mu_i}>0$ として
\begin{equation}\label{per:k4-critical-actions}
 I_{\mu_i}=-s,\qquad
 I_{\mu_j}=s\,\mathbf1_{\{j=f(i)\}}.
\end{equation}
graph-period correction は fiber tangent を消すので，この計算を変えない．

crossing face $ij$ の整数 torus basis は $(\mu_i,\mu_j)$ である．
$\epsilon_{ij}=\mathbf1_{\{j=f(i)\}}$ と書き，
\begin{equation}\label{per:k4-face-map}
 \begin{pmatrix}I_{\mu_i}\\ I_{\mu_j}\end{pmatrix}
 =A_{ij}\begin{pmatrix}t_i\\ t_j\end{pmatrix},\qquad
 A_{ij}=\begin{pmatrix}-1&\epsilon_{ji}\\
                         \epsilon_{ij}&-1\end{pmatrix},
 \qquad F=t_ir_i+t_jr_j
\end{equation}
で chart を定める．有向四角形には二辺 cycle がないので
$\det A_{ij}=1-\epsilon_{ij}\epsilon_{ji}=1$ である．
\eqref{per:k4-critical-actions} により，両端ではそれぞれ $F=sr_i$，$F=sr_j$ となる．

sector 内部で $t_i,t_j>0$ であることも確認する．
$i\to j$ が選択辺なら $t_i=-I_{\mu_i}$，
$t_j=-I_{\mu_i}-I_{\mu_j}$ である．$L_i$ 側の arm では，
partial basin area を用いた slope は $h_{ij}>-1$ なので
$t_j=-I_{\mu_i}(1+h_{ij})>0$ となる．$L_j$ 側では $h_{ji}>0$ なので，
やはり両係数は正である．いずれの向きも選択されない辺では
$F=-I_{\mu_i}r_i-I_{\mu_j}r_j$ で，両側の slopes は正である．
transition caps では radial coordinate を $I<w_j/2$ とし，normal moment $m_i$ を縮めれば
$I-w_j+m_i A(I)<0$ となる．product crossings でも二つの normal moments を
対応する weights より小さく取る．従って caps と corner の丸めを通じても上記の正値性を保つ．

作用座標は各 regular face で局所座標であり，Liouville flow はそれらを同じ倍率で変える．
従って contact-link edge 上の projective action ratio は，二つの端の rays の間で
厳密単調である．各 sector map は所期の fan sector の原点近傍への同相を与え，
\eqref{per:k4-critical-actions} によって隣接 sectors と貼り合う．
compact な link での一様な零極限が core への連続性を与える．
$A_{ij}$ は unimodular であり，$r_i,r_j$ もそれぞれの saturated rank-two lattice の
基底なので，face lattices と branch lattices は定義通り整合する．

最後に
\[
 H_1(W_4;\ZZ)=\ZZ^4/\ZZ(\mu_1+\mu_2+\mu_3+\mu_4)
\]
である\cite{Hattori,PapadimaSuciu}．
\eqref{per:k4-lattice-map} の四つの image の和は零であり，有向四角形の
連続する座標の差は係数和零の整数格子全体を生成する．
またその kernel は，cycle 全体で一定な係数の vectors，すなわち対角格子だけである．
従って $\Phi$ は整数格子同型である．
その $r_i,r_j$ 上の評価行列は正確に $A_{ij}$ なので，
\eqref{per:k4-face-map} は主張した linear-period 整合性を与える．
\end{proof}

この命題は一つの四直線 core に関する局所構成であり，
曲面全体の face transitions を含む atlas には，後の定理~\ref{ga8:global} の接着を用いる．
Evans--Mauri の contact-link と cone construction \cite{EM} に，
ここでは選んだ basin areas と作用格子の明示的な同定を加えている．



\section{ambient linear periods を要求した場合の分母二の障害}
\label{per:ambient-section}

前節の存在定理と矛盾しない，より強い同定への障害を述べる．
$N$ は標準 Petersen fan の rank-five tangent lattice，$H$ は
\eqref{per:global-meridians}の rank-four meridian group とする．
\emph{meridian-linear extension} とは homomorphism
\[
 \Phi:H\longrightarrow N^*
\]
であって，各 regular face 上の torus homology から $H$ への inclusion に対し，
$\Phi(\gamma)$ の restriction が前節の chart を介した実際の $I_\gamma$ に一致するものとする．
これは facewise action compatibility に追加する条件である．
affine functions を許しても，core における零極限が定数項を零にする．

\begin{theorem}[正規化された Petersen chart の ambient period defect]
\label{per:half-integral-obstruction}
定理\ref{per:local-realization}で構成した chart には，meridian-linear extension は存在しない．
値域を $N^*\otimes\mathbb Q$ に替えると extension は一意に存在し，
その各 meridian の値は半整数となる．
これらを含む最小の cotangent overlattice の $N^*$ に対する index は $16$ である．
この障害は追加した ambient extension に関するものであり，前節の unimodular face maps の障害ではない．
\end{theorem}

\begin{proof}
$\phi_a=\Phi(\mu_a)$ とする．crossing および matching faces 上の作用積分から
\[
 \phi_a(r_{A_a})=\phi_a(r_{B_a})=-1,
 \qquad
 \phi_a(r_{A_{a\pm1}})=\phi_a(r_{B_{a\pm2}})=0
\]
が必要である．標準 rays は
\[
 r_{A_{a-1}}+r_{A_{a+1}}+r_{B_a}=r_{A_a},
 \qquad
 r_{B_{a-2}}+r_{B_{a+2}}+r_{A_a}=r_{B_a}
\]
を満たす．これらを未定の四つの値に適用すると，一つの数 $x_a$ によって
\begin{equation}\label{per:ambient-values}
 \phi_a(r_{A_{a\pm2}})=x_a,
 \qquad
 \phi_a(r_{B_{a\pm1}})=1-x_a
\end{equation}
と書ける．逆に，この値の指定は全ての balancing relations を満たす．
$N=\operatorname{coker}(A-I)$ の presentation より，各 $x_a\in\mathbb Q$ は
一つの rational ambient linear form を定める．

meridian relation $\sum_a\mu_a=0$ は $\sum_a\phi_a=0$ を要求する．
これを $r_{A_j}$ 上で評価すると
\[
 x_{j-2}+x_{j+2}=1\qquad(j\in\ZZ/5\ZZ)
\]
となる．長さ五の奇周期に沿ってこの式を繰り返すと，$2x_a=1$ が全ての $a$ に対して従う．
$B_j$ 上の条件もこの解で満たされる．従って rational extension は一意で，
$x_a=1/2$ である．primitive ray 上で半整数値を取るため integral extension は存在しない．
matching face のもう一つの cycle $\beta_a=\mu_{a-1}+\mu_{a+1}$ も確認する．
$\phi_{a-1}+\phi_{a+1}$ は $r_{A_a}$ 上で零，$r_{B_a}$ 上で $1/2+1/2=1$ なので，
\eqref{per:internal-map} の $I_{\beta_a}$ と一致する．
従ってこの rational extension は個々の crossing meridians だけでなく，全ての face lattices 上で所要の作用積分を与える．

最後の index を直接計算する．命題\ref{trop:petersen-unimodular}と同じ基底
$[e_{12}],[e_{13}],[e_{14}],[e_{23}],[e_{24}]$ を使い，
$A_a\leftrightarrow\{2a,2a+1\}$，$B_a\leftrightarrow\{2a+2,2a+4\}$ とラベルを対応させる．
$2\phi_a$ の行列は
\begin{equation}\label{per:half-matrix}
 \begin{pmatrix}
 -2&1&0&1&0\\
 0&0&0&1&1\\
 1&0&1&0&-2\\
 1&1&-2&-2&1\\
 0&-2&1&0&0
 \end{pmatrix}.
\end{equation}
各列の和は零であり，modulo two でこの行列の rank は四である．
従って五つの $\phi_a$ の $\tfrac12N^*/N^*$ における像が生成する群は
$(\ZZ/2)^4$ であり，index は $2^4=16$ となる．
\end{proof}

\begin{remark}[homology diagram の向き]
\label{per:variance}
実際の branch block では $H_1(T_e)\to H_1(Q_v)$ は inclusion である．
一方，fan の cotangent lattices の自然な map は，star から face への restriction であり，向きが逆である．
従って $H_1(W_5)$ の rank 四と ambient lattice の rank 五を比較しただけでは障害は出ない．
定理\ref{per:half-integral-obstruction}は，追加条件 $\Phi$ を明示した上で，
実際の face restrictions と meridian relation を同時に解いた結果である．
整数座標で実現される各 face の作用積分と，それを ambient linear functions に延長する問題を区別する必要がある．
\end{remark}
